\input{Iteration-V2/Chapter5/chapter-05-12}

\subsection{Revision 30 evaluation: temporal correctness and bounded device evidence}

Revision 30 adds deterministic tests at the abstraction boundary where the temporal claim is made. Exact-boundary cases exercise timestamps immediately inside and outside a rule-owned window. A six-event fixture is evaluated under all 720 permutations to show that satisfying observations are order independent. A seeded 1,000-event burst checks bounded execution rather than presenting a synthetic timing result as a deployment benchmark. Additional cases verify profile validation, fail-closed rejection of invalid mappings, pack switching, source receipt generation, integration through the compliance engine, and wording that preserves notification-only behaviour. These tests supplement the retained 1,000-case engine corpus and 1,800-case independent governance oracle; they do not establish completeness against every Android signal or legal accuracy.

The release verification rebuilt the ARM64 release APK after the zero-frequency bridge guard and installed it on the recorded OPPO PERM00 device as version name 1.15.0 and version code 16. The APK SHA-256 was \texttt{\seqsplit{0bbaF15ba0a16379fae40c7d81f2ece09b20cb246a56939f46096fbb5d7ed2be}}. A cold launch reached the App, the UI tree exposed the Overview, Findings, and Settings controls, an on-device audit was requested, and the Findings evidence card was opened without a crash-buffer entry. This verifies a narrow executable path and readable local evidence workflow. It does not demonstrate a real-world temporal incident: the test device did not fabricate sensor activity solely to make an alert appear.

The evidential split is therefore intentional. Automated tests establish the exact temporal rule semantics, mapping substitution, fail-closed behaviour, and non-blocking integration; device inspection establishes installation and a representative user interaction path. Neither class of evidence proves GDPR compliance, universal Android coverage, legal interpretation, performance on all hardware, or user comprehension. The result is stronger traceability for a prototype, not a warrant for operational enforcement.
